% GENERATED FILE. Edit canonical lessons, then run npm run generate:books.
% canonical-source-hash: fnv1a64:069b7d679544ceda
% canonical-lessons: LA-C261-gigno, LA-C261-iubeo, LA-C261-cogo, LA-C261-lugeo, LA-C261-fodio

\chapter{To give birth to, To order, To force, To mourn, To dig}
\label{ch:la-to-give-birth-to-to-order-to-force-to-mourn-to-dig}
\hlchaptermodality{261}

\begin{chapteropening}
\textbf{By the end of this chapter:} \emph{I can say to give birth to, to order, to force, to mourn and to dig.}

Complete the last lesson of chapter 261: I can say to give birth to, to order, to force, to mourn and to dig.
\end{chapteropening}

\section[gignō, gignere]{gignō, gignere — to give birth to, to produce}

\label{lesson:LA-C261-gigno}



\begin{quote}
Before the new one: say the Latin for to notice, then the Latin for to blush.
\end{quote}

\subsection*{You'll want to know: gignō, gignere}
\textbf{gignō, gignere} — \textquotedblleft{}to give birth to, to produce\textquotedblright{}.

\begin{cousinweb}
\textbf{gignō} is inherited, usually traced to Proto-Indo-European \emph{*gi-gnh1-}, \textquotedblleft{}to beget\textquotedblright{}, a doubled form of the root (EtymDB 2.1); Greek \emph{gignomai}, \textquotedblleft{}to be born\textquotedblright{}, is a cousin (Lewis \& Short).
\end{cousinweb}

\begin{grammarlens}[title={What you've built}]
One more word for this chapter.
\end{grammarlens}

\subsection*{Your turn}
\begin{itemize}
\raggedright
  \item \emph{Say it:} \emph{gignō, gignere}
  \item \emph{Say it:} \emph{gignō, gignere}, once more
  \item \emph{Say it:} say \emph{gignō, gignere}
  \item \emph{Recall:} say the Latin for to notice, then the Latin for to blush, then say \emph{gignō, gignere} again
\end{itemize}

\begin{tcolorbox}[breakable,colback=teal!4,colframe=teal!35!black,title={Before you move on}]
What does gignō, gignere mean? (\textquotedblleft{}To give birth to, to produce\textquotedblright{}.) Say it once more.
\end{tcolorbox}
\section[iubeō, iubēre]{iubeō, iubēre — to order, to tell (someone to)}

\label{lesson:LA-C261-iubeo}



\begin{quote}
Before the new one: say the Latin for to blush, then the Latin for to give birth to.
\end{quote}

\subsection*{You'll want to know: iubeō, iubēre}
\textbf{iubeō, iubēre} — \textquotedblleft{}to order, to tell (someone to)\textquotedblright{}.

\begin{cousinweb}
\textbf{iubeō} is inherited; Old Latin wrote \emph{ioubeō}, and it is usually traced to Proto-Indo-European \emph{*Hyowdheye-}, \textquotedblleft{}to set in motion\textquotedblright{}, from \emph{*Hyewdh-}, \textquotedblleft{}to move\textquotedblright{} (EtymDB 2.1).
\end{cousinweb}

\begin{grammarlens}[title={What you've built}]
One more word for this chapter.
\end{grammarlens}

\subsection*{Your turn}
\begin{itemize}
\raggedright
  \item \emph{Say it:} \emph{iubeō, iubēre}
  \item \emph{Say it:} \emph{iubeō, iubēre}, once more
  \item \emph{Say it:} say \emph{iubeō, iubēre}
  \item \emph{Recall:} say the Latin for to blush, then the Latin for to give birth to, then say \emph{iubeō, iubēre} again
\end{itemize}

\begin{tcolorbox}[breakable,colback=teal!4,colframe=teal!35!black,title={Before you move on}]
What does iubeō, iubēre mean? (\textquotedblleft{}To order, to tell (someone to)\textquotedblright{}.) Say it once more.
\end{tcolorbox}
\section[cōgō, cōgere]{cōgō, cōgere — to force; to collect}

\label{lesson:LA-C261-cogo}



\begin{quote}
Before the new one: say the Latin for to give birth to, then the Latin for to order.
\end{quote}

\subsection*{You'll want to know: cōgō, cōgere}
\textbf{cōgō, cōgere} — \textquotedblleft{}to force; to collect\textquotedblright{}.

\begin{cousinweb}
\textbf{cōgō} is contracted from \emph{co-}, \textquotedblleft{}together\textquotedblright{}, and \emph{agō}, \textquotedblleft{}to drive\textquotedblright{} (Lewis \& Short): to drive together, to compel.
\end{cousinweb}

\begin{grammarlens}[title={What you've built}]
One more word for this chapter.
\end{grammarlens}

\subsection*{Your turn}
\begin{itemize}
\raggedright
  \item \emph{Say it:} \emph{cōgō, cōgere}
  \item \emph{Say it:} \emph{cōgō, cōgere}, once more
  \item \emph{Say it:} say \emph{cōgō, cōgere}
  \item \emph{Recall:} say the Latin for to give birth to, then the Latin for to order, then say \emph{cōgō, cōgere} again
\end{itemize}

\begin{tcolorbox}[breakable,colback=teal!4,colframe=teal!35!black,title={Before you move on}]
What does cōgō, cōgere mean? (\textquotedblleft{}To force; to collect\textquotedblright{}.) Say it once more.
\end{tcolorbox}
\section[lūgeō, lūgēre]{lūgeō, lūgēre — to mourn}

\label{lesson:LA-C261-lugeo}



\begin{quote}
Before the new one: say the Latin for to order, then the Latin for to force; to collect.
\end{quote}

\subsection*{You'll want to know: lūgeō, lūgēre}
\textbf{lūgeō, lūgēre} — \textquotedblleft{}to mourn\textquotedblright{}.

\begin{cousinweb}
\textbf{lūgeō} is probably inherited, from a Proto-Indo-European root \emph{*lewg-}, \textquotedblleft{}to break\textquotedblright{} (EtymDB 2.1); Greek \emph{lygros}, \textquotedblleft{}mournful\textquotedblright{}, is often compared (Lewis \& Short).
\end{cousinweb}

\begin{grammarlens}[title={What you've built}]
One more word for this chapter.
\end{grammarlens}

\subsection*{Your turn}
\begin{itemize}
\raggedright
  \item \emph{Say it:} \emph{lūgeō, lūgēre}
  \item \emph{Say it:} \emph{lūgeō, lūgēre}, once more
  \item \emph{Say it:} say \emph{lūgeō, lūgēre}
  \item \emph{Recall:} say the Latin for to order, then the Latin for to force; to collect, then say \emph{lūgeō, lūgēre} again
\end{itemize}

\begin{tcolorbox}[breakable,colback=teal!4,colframe=teal!35!black,title={Before you move on}]
What does lūgeō, lūgēre mean? (\textquotedblleft{}To mourn\textquotedblright{}.) Say it once more.
\end{tcolorbox}
\section[fodiō, fodere]{fodiō, fodere — to dig}

\label{lesson:LA-C261-fodio}



\begin{quote}
Before the new one: say the Latin for to force; to collect, then the Latin for to mourn.
\end{quote}

\subsection*{You'll want to know: fodiō, fodere}
\textbf{fodiō, fodere} — \textquotedblleft{}to dig\textquotedblright{}.

\begin{cousinweb}
\textbf{fodiō} is inherited, usually traced to Proto-Indo-European \emph{*bhod-}, \textquotedblleft{}to pierce, to dig\textquotedblright{} (EtymDB 2.1). English \emph{fossil}, \textquotedblleft{}dug up\textquotedblright{}, comes from its participle \emph{fossus}, \textquotedblleft{}dug\textquotedblright{}.
\end{cousinweb}

\begin{grammarlens}[title={What you've built}]
One more word for this chapter.
\end{grammarlens}

\subsection*{Your turn}
\begin{itemize}
\raggedright
  \item \emph{Say it:} \emph{fodiō, fodere}
  \item \emph{Say it:} \emph{fodiō, fodere}, once more
  \item \emph{Say it:} say \emph{fodiō, fodere}
  \item \emph{Recall:} say the Latin for to force; to collect, then the Latin for to mourn, then say \emph{fodiō, fodere} again
\end{itemize}

\begin{tcolorbox}[breakable,colback=teal!4,colframe=teal!35!black,title={Before you move on}]
What does fodiō, fodere mean? (\textquotedblleft{}To dig\textquotedblright{}.) Say it once more.
\end{tcolorbox}
